% ============================================================
% SEBC Dushigikirane — Description du Projet et Cadre
% Document technique de présentation
% ============================================================
\documentclass[12pt,a4paper]{article}
\usepackage[utf8]{inputenc}
\usepackage[T1]{fontenc}
\usepackage[french]{babel}
\usepackage{geometry}
\geometry{margin=2.5cm}
\usepackage{graphicx}
\usepackage{xcolor}
\usepackage{hyperref}
\usepackage{fancyhdr}
\usepackage{titlesec}
\usepackage{enumitem}
\usepackage{tabularx}
\usepackage{booktabs}
\usepackage{longtable}
\usepackage{tcolorbox}
\usepackage{fontawesome5}

% ═══ Couleurs SEBC ═══
\definecolor{sebcblue}{HTML}{1A3A5C}
\definecolor{sebcblulight}{HTML}{2C5F8A}
\definecolor{sebcred}{HTML}{C8102E}
\definecolor{sebcteal}{HTML}{0D9488}
\definecolor{sebcgray}{HTML}{64748B}
\definecolor{lightbg}{HTML}{F0F4FF}

% ═══ Styles ═══
\hypersetup{
    colorlinks=true,
    linkcolor=sebcblue,
    urlcolor=sebcblulight,
    citecolor=sebcblue,
}

\titleformat{\section}{\Large\bfseries\color{sebcblue}}{}{0em}{}[\titlerule]
\titleformat{\subsection}{\large\bfseries\color{sebcblulight}}{}{0em}{}
\titleformat{\subsubsection}{\normalsize\bfseries\color{sebcgray}}{}{0em}{}

\pagestyle{fancy}
\fancyhf{}
\fancyhead[L]{\small\color{sebcgray}SEBC Dushigikirane — Description du Projet}
\fancyhead[R]{\small\color{sebcgray}\thepage}
\fancyfoot[C]{\small\color{sebcgray}Document confidentiel — \today}
\renewcommand{\headrulewidth}{0.5pt}
\renewcommand{\footrulewidth}{0.3pt}

\tcbuselibrary{skins,breakable}

\newtcolorbox{infobox}[1][]{
    colback=lightbg,
    colframe=sebcblue,
    fonttitle=\bfseries,
    title=#1,
    rounded corners,
    boxrule=0.5pt,
    left=8pt, right=8pt, top=6pt, bottom=6pt,
}

\newtcolorbox{techbox}[1][]{
    colback=white,
    colframe=sebcteal,
    fonttitle=\bfseries\color{white},
    colbacktitle=sebcteal,
    title=#1,
    rounded corners,
    boxrule=0.5pt,
}

\begin{document}

% ═══════════════════════════════════════════════════════════
% PAGE DE TITRE
% ═══════════════════════════════════════════════════════════
\begin{titlepage}
\begin{center}
\vspace*{2cm}

{\Huge\bfseries\color{sebcblue} SEBC Dushigikirane}\\[0.5cm]
{\Large\color{sebcblulight} Soutien Entre Burundais du Canada}\\[2cm]

\rule{\textwidth}{2pt}\\[0.8cm]
{\LARGE\bfseries Description du Projet\\et Cadre Technique}\\[0.5cm]
\rule{\textwidth}{2pt}\\[2cm]

{\large\color{sebcgray}
Plateforme numérique de gestion associative\\
Messagerie interne · Vidéoconférence · Gestion des membres\\[1cm]
}

\begin{tabular}{ll}
\textbf{Version} & 1.0 \\
\textbf{Date} & \today \\
\textbf{URL Production} & \url{https://sebc-dushigikirane.pro} \\
\textbf{Développement} & Evanet08 \\
\end{tabular}

\vfill
{\small\color{sebcgray} Document confidentiel — Usage interne uniquement}
\end{center}
\end{titlepage}

% ═══════════════════════════════════════════════════════════
\tableofcontents
\newpage

% ═══════════════════════════════════════════════════════════
\section{Présentation Générale}
% ═══════════════════════════════════════════════════════════

\subsection{Contexte et Objectifs}

\textbf{SEBC Dushigikirane} (Soutien Entre Burundais du Canada) est une association à but non lucratif regroupant des membres de la diaspora burundaise résidant au Canada et dans d'autres pays. Le mot \textit{Dushigikirane} signifie « Entraidons-nous » en kirundi.

La plateforme numérique SEBC a été conçue pour répondre aux besoins suivants :

\begin{itemize}[leftmargin=2em]
    \item \textbf{Gestion centralisée des membres} : identification, adhésion, parrainage, approbation
    \item \textbf{Organisation structurelle} : cellules géographiques, types de membres hiérarchiques
    \item \textbf{Communication interne} : messagerie instantanée style WhatsApp, vidéoconférence intégrée
    \item \textbf{Gestion financière} : cotisations, soutiens en cas de décès, comptabilité
    \item \textbf{Administration modulaire} : système de modules dynamiques avec contrôle d'accès par rôle
\end{itemize}

\subsection{Portée du Projet}

\begin{infobox}[Périmètre fonctionnel actuel]
\begin{enumerate}[leftmargin=1.5em]
    \item Module \textbf{Authentification} : login par email, OTP, création de mot de passe
    \item Module \textbf{Candidature} : workflow en 2 étapes avec vérification du parrain
    \item Module \textbf{Mon Espace} (Membres) : tableau de bord personnel, gestion des ayants droits
    \item Module \textbf{Communication} : messagerie WhatsApp-style, vidéoconférence Jitsi, réunions planifiées
    \item Module \textbf{Administration} : gestion des paramètres, pays, provinces, cellules, modules, accès
    \item Module \textbf{Types de Membres} : hiérarchie Gestionnaire National / Cellule / Membre
\end{enumerate}
\end{infobox}

% ═══════════════════════════════════════════════════════════
\section{Architecture Technique}
% ═══════════════════════════════════════════════════════════

\subsection{Stack Technologique}

\begin{techbox}[Technologies utilisées]
\begin{tabularx}{\textwidth}{lX}
\toprule
\textbf{Composant} & \textbf{Technologie} \\
\midrule
Framework Backend & Django 6.0 (Python 3.10) \\
Base de données & MySQL / MariaDB \\
Serveur Web & Gunicorn (port 8004) + Apache Reverse Proxy \\
Hébergement & Serveur dédié Plesk (87.106.23.108) \\
Frontend & HTML5, CSS3 (design system custom), JavaScript vanilla \\
Emails transactionnels & API HTTP Brevo (SendinBlue) \\
Vidéoconférence & Jitsi Meet IFrame API (meet.jit.si) \\
Design System & Inter font, glassmorphism, pill gradient navigation \\
\bottomrule
\end{tabularx}
\end{techbox}

\subsection{Architecture Applicative}

L'application suit une architecture \textbf{monolithique Django} avec les couches suivantes :

\begin{enumerate}[leftmargin=2em]
    \item \textbf{Couche Modèle} (\texttt{models.py}) — 15+ modèles Django avec relations FK
    \item \textbf{Couche Vue} (\texttt{views.py}) — Pages rendues + APIs REST JSON (\texttt{@csrf\_exempt})
    \item \textbf{Couche Template} (\texttt{templates/}) — Templates Django avec héritage (\texttt{base.html})
    \item \textbf{Couche Service} (\texttt{email\_service.py}) — Services externes (Brevo API)
    \item \textbf{Couche Contexte} (\texttt{context\_processors.py}) — Injection dynamique sidebar
\end{enumerate}

\subsection{Modèle de Données}

Le schéma de base de données comprend les tables principales suivantes :

\begin{longtable}{p{4cm}p{9cm}}
\toprule
\textbf{Table} & \textbf{Description} \\
\midrule
\endhead
\texttt{membre} & Table centrale : identité, contact, cellule, rôle, type\_membre, auth \\
\texttt{type\_membre} & Hiérarchie des types : Gestionnaire (National/Cellule), Membre \\
\texttt{cellule} & Cellules géographiques (code, nom, province, pays) \\
\texttt{pays} & Référentiel pays (code ISO, indicatif téléphonique) \\
\texttt{province} & Provinces (rattachées à un pays) \\
\texttt{type\_ayant\_droit} & Types de liens familiaux \\
\texttt{ayant\_droit} & Bénéficiaires déclarés par le membre \\
\texttt{module} & Modules applicatifs (code, nom, icône, URL, ordre) \\
\texttt{acces\_module} & Matrice d'accès rôle × module (lire, écrire, supprimer) \\
\texttt{communication} & Messages de la messagerie interne (thread, scope, cibles) \\
\texttt{communication\_groupe} & Groupes de discussion personnalisés \\
\texttt{communication\_groupe\_membre} & Appartenance aux groupes \\
\texttt{meeting} & Réunions vidéo planifiées (room Jitsi, token join) \\
\texttt{meeting\_invite} & Invitations aux réunions \\
\texttt{notification\_gestionnaire} & Alertes d'administration \\
\texttt{parametre\_association} & Paramètres clé-valeur de l'association \\
\bottomrule
\end{longtable}

% ═══════════════════════════════════════════════════════════
\section{Système de Modules Dynamiques}
% ═══════════════════════════════════════════════════════════

\subsection{Principe}

Le système de navigation est entièrement \textbf{piloté par la base de données}. Chaque module applicatif est défini dans la table \texttt{module} avec :

\begin{itemize}
    \item \textbf{code} : identifiant unique (ex: \texttt{communication})
    \item \textbf{nom} : label affiché dans la sidebar
    \item \textbf{icone} : classe FontAwesome (ex: \texttt{fas fa-bullhorn})
    \item \textbf{couleur} : code hexadécimal pour l'avatar
    \item \textbf{url} : route Django associée
    \item \textbf{ordre} : position dans la sidebar
    \item \textbf{visible\_sidebar} : booléen de visibilité
    \item \textbf{est\_actif} : activation/désactivation globale
    \item \textbf{requiert\_approbation} : accès restreint aux membres approuvés
\end{itemize}

\subsection{Matrice d'Accès}

La table \texttt{acces\_module} associe chaque rôle à un module avec des permissions granulaires :

\begin{center}
\begin{tabular}{lccccc}
\toprule
\textbf{Module} & \textbf{Membre} & \textbf{Chef Cell.} & \textbf{Ch. Appro.} & \textbf{Comptable} & \textbf{Admin} \\
\midrule
Dashboard & — & R & R & R & RWD \\
Mon Espace & RW & RW & RW & RW & RWD \\
Communication & R & R & R & R & RWD \\
Administration & — & — & — & — & RWD \\
\bottomrule
\end{tabular}\\[0.3em]
{\small R = Lire, W = Écrire, D = Supprimer}
\end{center}

\subsection{Sidebar Dynamique}

La sidebar est générée par le \texttt{context\_processor} \texttt{sidebar\_modules} qui :
\begin{enumerate}
    \item Récupère le membre connecté depuis la session
    \item Filtre les modules actifs et visibles dans la sidebar
    \item Vérifie les droits d'accès selon le rôle du membre
    \item Injecte la liste filtrée dans le contexte de chaque page
\end{enumerate}

% ═══════════════════════════════════════════════════════════
\section{Hiérarchie des Types de Membres}
% ═══════════════════════════════════════════════════════════

La table \texttt{type\_membre} introduit une hiérarchie fonctionnelle :

\begin{infobox}[Types de membres]
\begin{tabularx}{\textwidth}{lllX}
\toprule
\textbf{Type} & \textbf{Niveau} & \textbf{Comm.} & \textbf{Description} \\
\midrule
Gestionnaire & National & \faCheck & Portée nationale, gestion complète \\
Gestionnaire & Cellule & \faCheck & Portée limitée à sa cellule \\
Membre & — & — & Membre régulier de l'association \\
\bottomrule
\end{tabularx}
\end{infobox}

Chaque type possède des flags de permissions :
\begin{itemize}
    \item \texttt{peut\_gerer\_communication} — Envoyer des messages de masse
    \item \texttt{peut\_approuver\_membres} — Approuver les candidatures
    \item \texttt{peut\_gerer\_finances} — Accès aux fonctions financières
\end{itemize}

% ═══════════════════════════════════════════════════════════
\section{Module Communication}
% ═══════════════════════════════════════════════════════════

\subsection{Architecture de Messagerie}

Le module Communication implémente une interface \textbf{WhatsApp-style} en plein écran avec :

\begin{itemize}
    \item \textbf{Sidebar gauche} : liste des contacts, cellules, groupes, recherche
    \item \textbf{Zone de chat droite} : messages avec séparateurs de dates, bulles envoyé/reçu
    \item \textbf{Barre d'envoi} : texte + pièces jointes (images, PDF, documents)
\end{itemize}

\subsubsection{Scopes de communication}

\begin{tabularx}{\textwidth}{lX}
\toprule
\textbf{Scope} & \textbf{Description} \\
\midrule
\texttt{individual} & Message privé entre deux membres \\
\texttt{cellule} & Message à tous les membres d'une cellule \\
\texttt{national} & Diffusion à tous les membres (gestionnaires uniquement) \\
\texttt{type\_membre} & Ciblage par type de membre \\
\texttt{custom\_group} & Groupe personnalisé créé par un membre \\
\bottomrule
\end{tabularx}

\subsubsection{Pièces jointes}

Les messages peuvent inclure des pièces jointes (max 10 MB) :
\begin{itemize}
    \item \textbf{Images} : JPG, PNG, GIF, WebP — affichage inline avec prévisualisation
    \item \textbf{Documents} : PDF, Word, Excel, PowerPoint — lien de téléchargement
    \item \textbf{Autres} : ZIP, RAR, TXT — téléchargement générique
\end{itemize}
Les fichiers sont stockés dans \texttt{media/communication/attachments/YYYY/MM/}.

\subsubsection{Système de threads}

Chaque conversation possède un \texttt{thread\_id} unique généré selon le scope :
\begin{itemize}
    \item \texttt{national} / \texttt{general} — Canaux globaux
    \item \texttt{cell\_\{id\}} — Fils de cellule
    \item \texttt{mbr\_\{min\}\_\{max\}} — Conversations privées (symétrique)
    \item \texttt{cgrp\_\{id\}} — Fils de groupe personnalisé
    \item \texttt{type\_\{id\}} — Fils par type de membre
\end{itemize}

\subsection{Vidéoconférence}

\begin{techbox}[Intégration Jitsi Meet]
La vidéoconférence est intégrée via l'\textbf{IFrame API de Jitsi Meet} :
\begin{itemize}
    \item \textbf{Fournisseur} : meet.jit.si (serveur public Jitsi, gratuit, open-source)
    \item \textbf{Participants} : nombre illimité
    \item \textbf{Sans publicité} : watermarks, promotions et bannières désactivés
    \item \textbf{Fonctionnalités} : audio, vidéo, partage d'écran, chat intégré
    \item \textbf{Configuration} : \texttt{SHOW\_JITSI\_WATERMARK: false}, \texttt{MOBILE\_APP\_PROMO: false}
\end{itemize}
\end{techbox}

Deux modes d'utilisation :
\begin{enumerate}
    \item \textbf{Appel direct} (\faVideo) — Génère un room éphémère et envoie automatiquement le lien dans le chat
    \item \textbf{Réunion planifiée} (\faCalendarPlus) — Création avec titre, date, durée, invitations + notifications email
\end{enumerate}

\subsection{Système de Réunions}

\begin{infobox}[Cycle de vie d'une réunion]
\begin{enumerate}
    \item \textbf{Planification} : l'organisateur crée la réunion (titre, date, durée, invités)
    \item \textbf{Notification} : email automatique envoyé aux invités via Brevo
    \item \textbf{Lien de partage} : URL unique avec token (\texttt{?join=TOKEN})
    \item \textbf{Rejoindre} : clic sur le lien → auto-join dans l'interface Jitsi embedded
    \item \textbf{Gestion} : liste des réunions (à venir / terminées), annulation possible
\end{enumerate}
\end{infobox}

Chaque réunion génère :
\begin{itemize}
    \item Un \textbf{room\_name} unique pour Jitsi (préfixé \texttt{SEBC-})
    \item Un \textbf{join\_token} UUID pour le lien de partage sécurisé
    \item Une \textbf{URL de partage} : \texttt{https://sebc-dushigikirane.pro/communication/?join=TOKEN}
\end{itemize}

\subsection{Messages non lus et notifications}

\begin{itemize}
    \item \textbf{Badge sur les contacts} : compteur rouge des messages non lus par thread
    \item \textbf{API dédiée} (\texttt{/api/communication/unread/}) : compteur total pour badge sidebar
    \item \textbf{Marquage automatique} : les messages sont marqués « lus » à l'ouverture du thread
\end{itemize}

% ═══════════════════════════════════════════════════════════
\section{Workflow d'Adhésion}
% ═══════════════════════════════════════════════════════════

Le processus d'adhésion se déroule en \textbf{deux étapes} :

\subsection{Étape 1 : Vérification du Parrain}
\begin{enumerate}
    \item Le candidat saisit l'email ou le téléphone du parrain
    \item Le système vérifie que le parrain existe dans la base
    \item Si trouvé, le candidat accède au formulaire d'identification
    \item Seul l'email du parrain est affiché (pas d'autres données)
\end{enumerate}

\subsection{Étape 2 : Identification et Vérification}
\begin{enumerate}
    \item Le candidat remplit ses informations personnelles complètes
    \item Un code OTP est envoyé par email (via Brevo API)
    \item Après vérification OTP, le candidat crée son mot de passe
    \item Le parrain reçoit une notification email pour valider le parrainage
    \item Option de relance du parrain si nécessaire
\end{enumerate}

% ═══════════════════════════════════════════════════════════
\section{APIs REST}
% ═══════════════════════════════════════════════════════════

L'application expose les endpoints suivants :

\begin{longtable}{p{6cm}p{2cm}p{5cm}}
\toprule
\textbf{Endpoint} & \textbf{Méthode} & \textbf{Description} \\
\midrule
\endhead
\multicolumn{3}{l}{\textbf{\color{sebcblue}Authentification}} \\
\texttt{/api/auth/check-email/} & POST & Vérifier existence email \\
\texttt{/api/auth/login/} & POST & Connexion \\
\texttt{/api/auth/request-otp/} & POST & Demander un OTP \\
\texttt{/api/auth/verify-otp/} & POST & Vérifier OTP \\
\texttt{/api/auth/set-password/} & POST & Créer mot de passe \\
\texttt{/api/auth/logout/} & POST & Déconnexion \\
\midrule
\multicolumn{3}{l}{\textbf{\color{sebcblue}Candidature}} \\
\texttt{/api/candidature/check-parrain/} & POST & Vérifier parrain \\
\texttt{/api/candidature/submit/} & POST & Soumettre candidature \\
\midrule
\multicolumn{3}{l}{\textbf{\color{sebcblue}Communication}} \\
\texttt{/api/communication/contacts/} & GET & Liste des contacts \\
\texttt{/api/communication/threads/} & GET & Threads avec derniers messages \\
\texttt{/api/communication/} & GET & Messages d'un thread \\
\texttt{/api/communication/send/} & POST & Envoyer un message (texte) \\
\texttt{/api/communication/send-file/} & POST & Envoyer avec pièce jointe \\
\texttt{/api/communication/unread/} & GET & Compteur messages non lus \\
\texttt{/api/communication/visio/} & GET & Générer room Jitsi \\
\texttt{/api/communication/groups/create/} & POST & Créer un groupe \\
\texttt{/api/communication/groups/delete/} & POST & Supprimer un groupe \\
\midrule
\multicolumn{3}{l}{\textbf{\color{sebcblue}Réunions}} \\
\texttt{/api/communication/meetings/} & GET & Liste des réunions \\
\texttt{/api/communication/meetings/create/} & POST & Planifier une réunion \\
\texttt{/api/communication/meetings/cancel/} & POST & Annuler une réunion \\
\texttt{/api/communication/meetings/join/} & GET & Rejoindre par token \\
\midrule
\multicolumn{3}{l}{\textbf{\color{sebcblue}Administration}} \\
\texttt{/api/admin/pays/} & POST & CRUD pays \\
\texttt{/api/admin/cellules/} & POST & CRUD cellules \\
\texttt{/api/admin/modules/} & POST & CRUD modules et accès \\
\bottomrule
\end{longtable}

% ═══════════════════════════════════════════════════════════
\section{Déploiement et Infrastructure}
% ═══════════════════════════════════════════════════════════

\begin{techbox}[Configuration de production]
\begin{tabularx}{\textwidth}{lX}
\toprule
\textbf{Paramètre} & \textbf{Valeur} \\
\midrule
Domaine & \texttt{sebc-dushigikirane.pro} \\
Serveur & Plesk Manager, IP : 87.106.23.108 \\
WSGI & Gunicorn, port 8004 \\
Proxy & Apache Reverse Proxy (HTTPS/SSL) \\
Base de données & MariaDB locale \\
Emails & Brevo API (émetteur : \texttt{link2ictgroup@gmail.com}) \\
Variables & \texttt{.env} : \texttt{BREVO\_API\_KEY}, \texttt{DEFAULT\_FROM\_EMAIL} \\
Service systemd & \texttt{sebc-dushigikirane.service} \\
\bottomrule
\end{tabularx}
\end{techbox}

\subsection{Workflow de déploiement}

\begin{enumerate}
    \item Développement local → \texttt{makemigrations}
    \item \texttt{rsync} des fichiers vers le serveur de production
    \item \texttt{python manage.py migrate} sur le serveur
    \item \texttt{python init\_data.py} pour les données de référence
    \item \texttt{systemctl restart sebc-dushigikirane}
\end{enumerate}

% ═══════════════════════════════════════════════════════════
\section{Sécurité}
% ═══════════════════════════════════════════════════════════

\begin{itemize}
    \item \textbf{Authentification} : hash SHA-256 avec sel aléatoire (16 caractères)
    \item \textbf{Sessions} : session Django standard avec vérification par \texttt{\_get\_logged\_membre}
    \item \textbf{OTP} : codes à 6 chiffres avec expiration (15 minutes)
    \item \textbf{CSRF} : désactivé sur les APIs JSON (\texttt{@csrf\_exempt}) — à sécuriser en production
    \item \textbf{Uploads} : limite 10 MB, types de fichiers restreints
    \item \textbf{Meetings} : tokens UUID uniques pour les liens de join
\end{itemize}

% ═══════════════════════════════════════════════════════════
\section{Évolutions Prévues}
% ═══════════════════════════════════════════════════════════

\begin{enumerate}
    \item Module \textbf{Cotisations} : suivi des paiements, historique, rappels automatiques
    \item Module \textbf{Soutien} : gestion des cas de décès, répartition des fonds
    \item Module \textbf{Rapports} : tableaux de bord statistiques, exports PDF
    \item \textbf{Notifications push} : intégration Web Push pour les nouveaux messages
    \item \textbf{Application mobile} : PWA ou application native React Native
    \item \textbf{Serveur Jitsi dédié} : hébergement propre pour la vidéoconférence
\end{enumerate}

\vfill
\begin{center}
\rule{0.5\textwidth}{0.5pt}\\[0.5em]
{\small\color{sebcgray} SEBC Dushigikirane — version 1.0 — \today}
\end{center}

\end{document}
